% =====================================================================
%  sec-metalanguage-new.tex
%
%  Drop-in section for Jon's paper `algebraic-elaboration'.
%  DESTINATION: sourcehut/algebraic-elaboration/main.tex
%  PLACEMENT:   paste as a NEW SECTION immediately BEFORE
%                   \section{Modular elaborators in the monadic metalanguage}
%               (currently ~line 592 of that file), i.e. directly after the
%               \mathbf{scope}-notation subsection of \S\ref{sec:metalanguage}.
%               It is meant to subsume the unfinished "Compact types ... scope"
%               material and the \varpi sketch in trash.tex.
%
%  PREAMBLE DEPENDENCIES (all already present in the paper):
%     tikz-cd ; amsmath ; the theorem environments `definition' and `remark'.
%  The reference \ref{con:lax-monoidal-compact} points at the existing
%  Construction in the paper. The citation altenkirch-chapman-uustalu:2015
%  must be added to refs.bib (bibtex given in the accompanying chat message).
%
%  TO PREVIEW inside your own top-level main.tex instead, add to its preamble:
%     \newtheorem{remark}{Remark}
%  and either delete the \ref/\citet or provide stubs. Everything else uses
%  only standard LaTeX (\prod, \to, \jmath, tikz-cd).
% =====================================================================



\begin{remark}[Reduction to the ordinary case]
When $J=\mathrm{Id}_{\mathcal{C}}$ (so $\mathcal{C}=\mathbf{Set}$,
$\jmath=\mathrm{id}$, and $g^{*}=\mu\circ Tg$), the data is an ordinary
$\mathcal{K}$-ary lax monoidal monad and axioms (4)--(5) reduce to its unit-
and multiplication-compatibility. For (5): taking $g_k=\mathrm{id}_{TY_k}$ (so
$g_k^{*}=\mu_{Y_k}$, $\widehat{g}=\phi_{K,Y}$ and
$(\widehat g)^{*}=\mu\circ T\phi_{K,Y}$) turns the bind square into
\[
  \phi_{K,Y}\circ\textstyle\prod_k\mu_{Y_k}
  \;=\;
  \mu_{\prod Y}\circ T\phi_{K,Y}\circ\phi_{K,TY},
\]
which is exactly multiplication-compatibility; conversely, that axiom together
with naturality of $\phi$ reproduces the bind square for an arbitrary family
$g$. Passing to $(-)^{*}$ is the \emph{only} change forced by the typing
$T\colon\mathbf{Set}\to\mathcal{C}$: the non-relative term $\phi_{K,TX}$ would
require $T$ to be applied to the objects $TX_k\in\mathcal{C}$, which is
ill-typed for a relative monad. In short, ``multiplication is monoidal''
becomes ``Kleisli extension is monoidal''.
\end{remark}

\begin{remark}[Why surjections, and not arbitrary maps: weakening]
\label{rem:weakening}
Every re-indexing $f\colon K\to L$ factors as a surjection followed by an
injection. The surjective part duplicates indices --- a \emph{diagonal} on
products --- and corresponds logically to \emph{contraction} (using one bound
value in several positions); the injective part deletes indices --- a
\emph{projection} on products --- and corresponds to \emph{weakening}
(discarding an unused hypothesis). Compatibility with the projection part is
equivalent to $T$ preserving the corresponding products, and \emph{fails} for
the partiality monad: already for the lift monad on $\mathbf{Set}$ one has
$T(\pi_1)\circ\phi_{2}\neq\pi_1$, because $\phi_2$ has conjoined the definedness
of the very coordinate that $\pi_1$ discards. In the internal language,
$\mathbf{scope}$-ing a variable $k\colon K$ that is not used still imposes
$\forall k\colon K$ on the resulting definedness, which is strictly stronger
than not binding it at all. We therefore require re-indexing only along
surjections; this still subsumes the symmetric-group action (renaming of bound
variables) and contraction, both of which the partiality monad validates.
Should a particular application genuinely require weakening through
$\mathbf{scope}$, it must be imposed as a separate axiom, and doing so excludes
the partiality monad.
\end{remark}

\begin{remark}[The concrete instance]
Take $\mathcal{C}$ to be the ambient topos (or $\mathbf{Pos}$), $J$ the discrete
embedding, and $T=\mathbf{L}_\Sigma$ the partiality relative monad. The family
$\phi_{K,X}$ is then the structure $\boldsymbol{\vartheta}_{K,X}$ of
Construction~\ref{con:lax-monoidal-compact}, which exists exactly when the arity
$K$ is $\Sigma$-compact (so that $\forall k\colon K$ preserves
$\Sigma$-openness). Under this dictionary, axiom (4) is the unit equation of the
$\mathbf{scope}$-notation, axiom (5) is its bind equation, and axiom (3) holds
for surjections. Thus the abstract definition specialises exactly to the
combinator already used to interpret binders.
\end{remark}

\begin{remark}[Abstract form]
Equivalently, in the manner in which a $J$-relative monad is a monoid in the
skew-monoidal category $([\mathbf{Set},\mathcal{C}],\diamond_J,J)$ with
$F\diamond_J G=\mathrm{Lan}_J F\circ G$, a $\mathcal{K}$-ary lax monoidal
structure makes $T$ a monoid for $\diamond_J$ \emph{inside} the category of
$\mathcal{K}$-ary lax monoidal functors $\mathbf{Set}\to\mathcal{C}$ --- i.e.\
both the unit $J$ and the tensor $\diamond_J$ are required to lift to that
category. We do not use this packaging below; it is recorded only to locate the
definition in the literature on substitution-monoidal structures.
\end{remark}
